\chapter{The Formal Seed}\label{ch11:lean}

Chapter~10 documented the executable evidence: bounded-exhaustive suites
that tie the running Python to the proved mathematics wherever a finite
check can reach. This chapter documents the artifact's second, independent
evidence channel: a small body of machine-checked mathematics --- the
\emph{formal seed} --- developed in Lean~4 under a pinned toolchain and
audited by a dedicated verifier script. The seed is deliberately modest in
scope and immodest in rigor: what it checks, it checks down to the proof
kernel with a single standard axiom; what it does not check, this chapter
itemizes without euphemism. The governing records are the manifest's
\texttt{formal\_verification} block (status \textsc{partial / present}),
the integrated Lean receipt \texttt{LEAN\_PROOF\_MANIFEST\_v0.1.json}, and
the lift-functor subreceipt, all hash-bound alongside the code they
describe.

\section{Kernel-checking as a trust story}\label{ch11:sec:trust}

\subsection{Why a proof assistant, and what its word is worth}

An informal proof --- even a careful one, even one checked by several
readers --- is a social artifact: its correctness claim rests on the
attention of the humans who have read it. A proof developed in Lean~4
\cite{lean4} rests on a different foundation. Lean is architected around a
small trusted kernel: the elaborate machinery that makes proofs writable
--- tactics, automation, syntax extensions, metaprogramming --- all runs
\emph{outside} the trust boundary, producing proof terms that the kernel
independently re-checks against the rules of the underlying dependent type
theory. To trust a Lean-checked theorem one must trust the kernel's
implementation of a well-studied type theory, the statement as written,
and the axioms the proof depends on --- and nothing else. Not the tactics
that found the proof, not the libraries that suggested it, not the human
who wrote it.

This architecture is the modern form of what the proof-assistant
literature calls the de Bruijn criterion: the system's trustworthiness
should rest on a small, independently auditable checker, with all
proof-\emph{finding} machinery outside the trusted base. Lean~4
satisfies the criterion by construction --- the kernel implements a
typechecker for its dependent type theory and nothing else, and every
theorem, however it was found, is ultimately a term the kernel has
typechecked against its stated proposition. The practical consequence
for this artifact is a sharp asymmetry of failure modes. If a tactic is
buggy, the worst case is a proof that fails to build --- annoying,
visible, safe. For a false statement to be \emph{accepted}, the defect
must be in the kernel itself, in the type theory it implements, or in
the statement of the theorem --- and the first two are shared,
heavily exercised infrastructure whose failure would be a community
event, not a private one. The third --- a theorem statement that does
not say what its author believes it says --- is the one failure mode
kernel-checking cannot address at all, and it is why the verifier
discipline below insists on hash-bound sources and named theorem
inventories: the statements themselves must be reviewable text, pinned
to digests, so that what was checked is exactly identifiable.

That last dependency --- the axioms --- is where formal verification
claims most often quietly leak. Lean permits axiom declarations, and its
standard library itself provides classical axioms; a development can also
locally postulate whatever it likes, and a proof ``checked by Lean'' may
silently rest on an unproved assumption three imports away. Similarly, the
placeholder \texttt{sorry} lets a development compile with holes. A
verification claim is therefore only as strong as its \emph{axiom audit}:
an explicit accounting, per theorem, of exactly which axioms the checked
proof term depends on, together with a policy on which axioms are
admissible. The seed's verifier makes this audit mechanical, as
Section~\ref{ch11:sec:verifier} details: the sole admissible standard
axiom is \texttt{propext} --- propositional extensionality, the principle
that logically equivalent propositions are equal, a standard and
deliberately minimal extension of the core theory --- and any project
axiom or \texttt{sorry} anywhere in the build fails the verification run
outright.

\subsection{Why \texttt{propext}, and why only \texttt{propext}}

The single admitted axiom deserves its own paragraph, because the choice
encodes a policy. Propositional extensionality states that two
propositions which imply each other are equal as propositions --- an
identification of logical equivalence with equality at the level of
\texttt{Prop}. It is among the mildest of the standard classical
extensions: it does not decide any proposition (it is not excluded
middle), it does not choose witnesses (it is not choice), and it is
validated by the standard proof-irrelevant semantics of \texttt{Prop}.
Lean's own core library depends on it pervasively, which makes it the
practical floor for any development that uses the standard library at
all; a seed that admitted nothing would have to rebuild elementary
logic from scratch for no gain in trust. The policy, then, is: admit
the one axiom that the ecosystem's floor requires, and nothing above
that floor. In particular the seed does not admit choice, quotient
soundness beyond the kernel's built-in support, or classical excluded
middle --- and because the verifier audits dependencies per theorem,
the claim is not ``we tried to avoid stronger axioms'' but ``the
compiled proof terms of the named theorems depend on \texttt{propext}
and nothing else, mechanically checked.'' For the finite, discrete,
computation-adjacent statements in the verified scope, this is also
mathematically comfortable: nothing in typed identities,
associativity, basis compositions, or two-state expansions has any
classical content to hide.

\subsection{The seed's place in the evidence architecture}

The formal seed occupies a precise position between the other two evidence
channels. Against the informal proofs of Chapters~4--8, it trades
coverage for certainty: only a fraction of the monograph's mathematics is
formalized, but that fraction is checked by a mechanism whose failure
modes are qualitatively narrower than human review. Against the
executable suites of Chapter~10, it trades concreteness for universality:
the suites exercise the actual Python on bounded instances, while the
seed proves \emph{universal} statements --- for all mappings, all
carriers of the modeled form --- about a Lean model that is not the
Python. Neither channel subsumes the other, and the artifact's claim
discipline insists on reporting them separately: the manifest's
formal-verification status is \textsc{partial / present}, its claim field
stating that the seed ``kernel-checks only the exact modules and theorem
dependencies named in the linked proof manifest,'' with everything else
\textsc{open}.

The three channels can be tabulated by what each quantifies over and
what each trusts:

\begin{center}
\small
\begin{tabular}{@{}lp{3.4cm}p{3.4cm}p{3.6cm}@{}}
\toprule
Channel & Quantifies over & Object of evidence & Trusted base \\
\midrule
Informal proof & all instances & the mathematics & human review of the argument \\
Executable suite & bounded instances & the running Python & interpreter, test oracles, bound choice \\
Lean seed & all instances (in scope) & a Lean model & Lean kernel, type theory, statement fidelity \\
\bottomrule
\end{tabular}
\end{center}

The table makes the complementarity mechanical. The suite is the only
channel that touches the shipped implementation; the seed is the only
channel whose checking is not human; the informal proofs are the only
channel with full coverage. Any two channels overlapping on a statement
--- as all three do on the noncontractivity fixture --- cross-check
each other's distinct trusted bases, which is the strongest evidential
configuration the artifact attains and the reason the fixture is
maintained in all three forms.

\section{The pinned toolchain}\label{ch11:sec:toolchain}

Formal verification inherits the reproducibility obligations of any other
evidence. A proof that checks under one compiler version and fails under
another --- because of changed elaboration, moved library lemmas, or a
revised axiom inventory --- is not a stable receipt. The seed therefore
pins its entire toolchain, and the manifest's environment probe of
2026-08-04 records the pin exactly:

\begin{center}
\begin{tabular}{@{}ll@{}}
\toprule
Component & Recorded state \\
\midrule
Lean & 4.32.1 / \texttt{f054605aea4b840552cca2e725580bffd1e1b704} / \textsc{present} \\
Lake & 5 / \textsc{present} \\
Python (verifier host) & 3.14.0 \\
\texttt{coqc} & \textsc{unavailable} \\
Agda & \textsc{unavailable} \\
\bottomrule
\end{tabular}
\end{center}

The Lean entry pins not only the release number (4.32.1) but the exact
toolchain commit hash, so ``the toolchain'' is an identified artifact
rather than a version-string family. The build goes through Lake, Lean's
standard build system, so module dependency resolution is itself
declarative and replayable. The toolchain is furthermore \emph{isolated}
--- pinned for this project rather than taken from whatever ambient
Lean installation the host happens to carry --- so a developer machine
with a different global Lean cannot silently substitute its own
elaboration behavior into a verification run. The recorded
\textsc{unavailable} entries for Coq and Agda are honesty about the
probe environment, with a methodological implication worth noting: the
seed is a single-prover development, and cross-prover corroboration ---
re-checking the same statements in an independent system with an
independent kernel --- is one of the things the artifact does not have.

Enumerating the trusted computing base explicitly, in decreasing order
of scrutiny received: the Lean kernel at the pinned commit; the type
theory it implements; the Lake build tool's module resolution; the
Python 3.14.0 interpreter running the verifier script; the verifier
script itself --- which is why the script is kept small, its policy
constants hardcoded rather than configurable, and its own bytes
hash-bound in the artifact envelope; and the host operating system
underneath all of it. The pin does not shrink this base; it makes the
base \emph{identified}, which is what turns ``it verified'' from an
anecdote about some machine into a replayable claim about a named
configuration.

\section{The verifier discipline}\label{ch11:sec:verifier}

The script \texttt{scripts/verify\_psr8\_lean\_gate1.py} is to the formal
seed what the Gate~1 verifier is to the executable suite: a single
replayable command that either reproduces the full evidence or fails.
Its checks, as recorded in the specification and the research paper,
are these:

\begin{protocol}[Lean seed verification]\label{ch11:proto:lean}
The verifier, run as
\texttt{python scripts/verify\_psr8\_lean\_gate1.py}:
\begin{enumerate}
\item \emph{hardcodes the sole allowed standard axiom},
      \texttt{propext} --- the admissible-axiom set is not a
      configuration option but a constant in the verifier source;
\item \emph{rejects project axiom declarations} --- any
      \texttt{axiom} introduced by the development itself fails the run,
      so the seed cannot smuggle assumptions in as local postulates;
\item \emph{rejects \texttt{sorry}} --- no proof hole anywhere in the
      built modules is tolerated;
\item \emph{checks source hashes} --- the Lean source files are compared
      by digest against the proof manifest, so the theorems checked are
      the theorems bound, not edited neighbors;
\item \emph{builds through Lake} under the pinned toolchain of
      Section~\ref{ch11:sec:toolchain}; and
\item \emph{audits exact theorem dependency reports} --- for each named
      theorem, the axioms its compiled proof term actually depends on are
      extracted and compared against the expected inventory, so the audit
      is per-theorem and mechanical, not a global grep.
\end{enumerate}
\end{protocol}

The design mirrors, in the formal channel, the substitution-closure of the
Gate~1 protocol: hash checks prevent verifying edited sources; the axiom
and \texttt{sorry} rejections prevent weakening the logic; the dependency
audit prevents an individually honest theorem from silently importing a
dishonest helper. What survives all six checks is a set of named
statements, in identified source files, proved to the kernel's
satisfaction from \texttt{propext} alone under an identified compiler.

\begin{remark}[Scope of the receipt]\label{ch11:rem:receipt}
The integrated receipt \texttt{LEAN\_PROOF\_MANIFEST\_v0.1.json}
(hash-bound in the prototype manifest, with the lift-functor subreceipt
\texttt{LEAN\_DETERMINISTIC\_LIFT\_FUNCTOR\_MANIFEST\_v0.1.json} bound
beside it) names the exact modules and theorem identifiers in scope. The
prototype manifest's verification claim defers to that receipt ---
``only the exact modules and theorem dependencies named in the linked
proof manifest'' --- and this chapter's itemization follows the same
inventory as summarized in the manifest's \texttt{verified\_scope} and
\texttt{open\_scope} fields.
\end{remark}

\section{Verified scope}\label{ch11:sec:verified}

The seed's kernel-checked content falls in seven items, quoted here from
the manifest's verified-scope inventory and expanded with the
mathematical identity of each.

\begin{enumerate}
\item \emph{Typed identities and associativity, with endomap-only
      powers.} The Lean module \texttt{TypedCategory} models the kernel's
      typed total-function discipline and proves the identity laws and
      associativity of composition --- the formal counterpart of the
      category laws of Chapter~4 --- together with the powering
      discipline that admits powers only for endomaps, the formal
      counterpart of the endomap/path separation of Chapters~5 and~6.
      This is the seed's contribution to PO-TKS8-CORE-01, whose ledger
      status \textsc{partial formal verification present} rests on
      exactly these lemmas.
\item \emph{The exact two-state constant-lift expansion.} For the
      constant map on a two-point carrier --- the permanent
      counterexample of Chapter~7, ledger record
      \textsc{thm-lift-occupation-counterexample} --- the seed
      kernel-checks the exact vector expansion of the lift's action and
      its \emph{noncontractivity}: the checked statements exhibit the
      concrete vector on which the constant lift strictly increases
      norm, blocking forever the general claim that deterministic lifts
      are contractive. The fixture thus exists in all three evidence
      channels at once: informal proof, executable regression, and
      kernel-checked lemma.
\item \emph{Lift identity, entrywise and on basis vectors.} The
      statements $A_{\id_X} = I$ are checked in two forms: entrywise
      equality of the modeled matrix, and action on standard basis
      vectors.
\item \emph{Typed basis composition.} The composition half of the
      functor law is checked at basis level:
      $A_g(A_f e_x) = e_{g(f(x))} = A_{g\circ f} e_x$ for the modeled
      typed maps --- the on-basis-vectors content of
      Chapter~7's functoriality theorem.
\item \emph{Relational-product support composition.} At the level of
      supports --- which target coordinates can be nonzero --- the seed
      checks that the support of a composite lift is the relational
      product of the component supports, tying the matrix picture to the
      underlying relation composition.
\item \emph{Four-atom finite-context lemmas.} The seed includes the
      exhaustiveness and witness-map lemmas for the four-atom ACBE/RPM
      finite context used by the companion Noetic Topos development ---
      finite case analyses whose value is precisely that a kernel, not a
      human, confirmed the case split is total. Their presence in this
      seed is an artifact of shared toolchain and shared verifier; their
      mathematical home is the companion paper, and no TKS/CAS semantic
      claim travels with them (the boundary of
      Section~\ref{ch11:sec:nonverification} applies with full force).
\end{enumerate}

Together with the endomap-only power discipline these items constitute
the entirety of PO-TKS8-CORE-07's \textsc{partial / present} status and
the formal fraction of PO-TKS8-CORE-11 (\textsc{partial formal
verification present / finite-sum matrix bridge open}).

\section{Open scope}\label{ch11:sec:open}

The seed's honesty is its itemized deficit. The following are
\emph{not} formalized; each entry names the obligation it blocks or
partially blocks.

\begin{enumerate}
\item \emph{The universal norm theorem.} The exact operator-norm formula
      $\lVert A_f\rVert^2 = \max_y |f^{-1}(y)|$ (Chapter~7) is checked
      in Lean only for the two-state constant expansion. The universal
      statement --- all carriers, all mappings --- has an informal
      fiber-block Gram proof and bounded-exhaustive structural evidence,
      but no kernel-checked form. (PO-TKS8-CORE-05: \textsc{informal
      proof present / universal formalization open}.)
\item \emph{The generic matrix bridge.} The checked two-state lemmas
      speak of a concrete expansion, not of the general definitions.
      The bridge from the checked \texttt{constantTwoLift} object to
      generic matrix, Gram, trace, and rank definitions --- the step
      that would make the counterexample a corollary of a general
      theory rather than a standalone computation --- is open.
\item \emph{Finite-sum matrix multiplication.} The seed does not
      formalize numerical matrix multiplication as a finite sum, which
      is the missing link between the basis-level functor law it checks
      and the entrywise matrix functor law the Python tests check.
      This is the ``finite-sum matrix bridge'' of PO-TKS8-CORE-11.
\item \emph{Arbitrary-vector linearity.} The checked lift laws act on
      basis vectors; the linear extension to arbitrary vectors ---
      routine on paper, but a genuine formalization step --- is open.
\item \emph{The factorial idempotent-power theorem.} Theorem
      $f^{|X|!} \circ f^{|X|!} = f^{|X|!}$ (Chapter~5) is a standard
      result with a complete informal proof and a 289-case exhaustive
      check, and no Lean form. (PO-TKS8-CORE-04.)
\item \emph{Shortlex leastness.} The soundness of representative-first
      pruning and the leastness of the computed normal forms
      (Chapter~8) rest on the informal induction of the pruning lemma.
      (PO-TKS8-CORE-06.)
\item \emph{Presheaf, sheaf, topology-generation, and product-site
      layers.} The full site-theoretic program of the companion Noetic
      Topos development --- everything beyond the four-atom finite
      context lemmas --- is unformalized.
\end{enumerate}

The pattern across the list is consistent: the seed formalizes
\emph{discrete, finite, definition-adjacent} statements --- laws that
are within short reach of the modeled definitions --- and leaves open
everything requiring real analytic or combinatorial infrastructure
(spectral bounds, factorials and cycle structure, well-founded
minimization arguments, sheaf theory). This is a defensible frontier
for a seed, and the word ``seed'' is doing exact work: the artifact
claims a viable, pinned, audited starting point, not a formalized
theory.

\subsection{Formalization routes for the open items}

Without promising schedules --- the ledger records obligations, not
plans --- the open items admit an informed difficulty ordering, and
recording it here serves whoever next extends the seed.

The \emph{finite-sum matrix bridge} and \emph{arbitrary-vector
linearity} are the mechanical end: both are finite summation
arguments over Kronecker deltas, requiring only finite-sum
manipulation lemmas (reindexing, splitting over a fiber
decomposition), and both would immediately upgrade PO-01-adjacent and
PO-11 evidence by connecting the checked basis-level laws to the
entrywise matrix laws the Python suite tests. They are the natural
first targets.

The \emph{factorial power theorem} is middle difficulty: its proof
needs a formalized functional-graph decomposition (every point's
orbit is eventually periodic; tail length bounded by $|X| - 1$) and
the divisibility fact $\ell \le |X| \Rightarrow \ell \mid |X|!$. Both
components are within reach of standard library support for
\texttt{Nat} arithmetic and finite iteration, and an orbit-based
proof avoiding explicit graph theory --- iterate the pigeonhole
argument on the sequence $x, f(x), f^2(x), \dots$ --- is likely the
short path.

The \emph{universal norm theorem} has two routes with sharply
different formalization costs. The spectral route of the informal
proof (Gram block decomposition, largest eigenvalue) imports spectral
theory of symmetric matrices --- heavy machinery for a finite
statement. The elementary route bounds
$\lVert A_f v \rVert^2 = \sum_y \bigl|\sum_{x \in f^{-1}(y)} v_x\bigr|^2$
by Cauchy--Schwarz within each fiber, obtaining
$\le \max_y |f^{-1}(y)| \cdot \lVert v \rVert^2$, and exhibits the
normalized indicator of a maximal fiber as the equality witness ---
finite sums and one inequality, no spectra. The elementary route
depends only on the linearity item above and is plainly the
formalization-friendly path.

\emph{Shortlex leastness} is the most program-like: it requires an
abstract model of the closure exploration (a well-founded induction
over shortlex order with the pruning invariant) and is best attacked
as algorithm verification over an idealized queue rather than as a
lemma about the Python. The \emph{presheaf and sheaf layers} are the
far end --- genuine site theory --- and belong to the companion
program's formalization lane rather than the kernel's.

\section{What kernel-checking does not verify}\label{ch11:sec:nonverification}

\begin{boundary}[Model theorems are not program verification]\label{ch11:bdy:model}
Every theorem in the seed is a statement about \emph{Lean objects}: a
Lean model of typed finite functions, a Lean expansion of a two-state
lift. No theorem in the seed mentions, models, or constrains the Python
source of \texttt{src/tks\_v8\_core}. The relationship between the Lean
model and the running implementation is exactly the relationship the
red-team methodology treats it as: an informal correspondence, argued
by inspection --- the definitions were written to match --- and
checked on bounded instances by the executable suite, but never itself
proved. A reader who concludes ``the kernel is formally verified''
has crossed a boundary that the artifact's own manifest refuses to
cross: the manifest's claim records keep ``equivalence between the
Python implementation and the Lean model'' explicitly \textsc{open}.
\end{boundary}

The gap is not pedantry; it decomposes into three concrete failure
channels, none excluded by the seed. First, \emph{modeling error}: the
Lean definitions could mismatch the Python semantics --- an off-by-one
in the modeled indexing, a subtly different totality convention --- in
which case the checked theorems are true of the model and silent about
the code. Second, \emph{coverage error}: the Python contains behavior
the model omits entirely (resource guards, error taxonomy, presentation
naming, the closure engine), so even a perfect model of the composition
core would verify a fraction of the artifact. Third, \emph{compilation
and runtime error}: the Python executes on an interpreter whose own
correctness is assumed, not proved --- the standard trusted-computing
caveat that verified-compilation projects address and this artifact
does not attempt. Programs proved correct in the strong sense require
either verification-aware extraction (the program is generated from the
checked development) or a verified semantics of the implementation
language; the seed does neither, claims neither, and the monograph
records the fact in the ledger rather than in a footnote.

What the seed \emph{does} deliver, within these limits, is
significant: the mathematical laws at the heart of the design --- the
category laws, the power discipline, the counterexample that anchors
the norm theory's claim boundary --- hold in at least one fully
precise, kernel-audited model of the definitions, under a single
standard axiom, reproducibly, on a pinned toolchain. Errors in the
\emph{mathematics} of Chapters~4--7, if any survived informal review,
would have had to survive the kernel too wherever the scopes overlap.
The seed thus hardens the theory; it leaves the implementation's
verification status exactly where Chapter~10 put it.

\section{The PO-07 partial status}\label{ch11:sec:po07}

Obligation PO-TKS8-CORE-07 --- ``independent theorem-prover
formalization'' --- carries the ledger status \textsc{partial /
present}, and this chapter is that status rendered in full. It is
\textsc{present}: a pinned Lean~4.32.1 seed exists, builds through
Lake, passes a hardcoded single-axiom audit with project axioms and
\texttt{sorry} rejected, has its sources hash-bound, and checks the
seven-item scope of Section~\ref{ch11:sec:verified}. It is
\textsc{partial}: the seven-item open scope of
Section~\ref{ch11:sec:open} remains, the Python-to-model equivalence
is unproved, and no second prover corroborates the first. The word
``independent'' in the obligation's title is also only partially
honored in a further sense worth recording: the seed was developed
within the project, so while the \emph{checking} is independent of
human attention in the way only a kernel can be, the \emph{statement
selection} is not independent of the authors --- a reviewer could ask
whether the formalized statements are the load-bearing ones, and
PO-TKS8-CORE-10 (independent specialist review, \textsc{open}) is
where that question lives. Closing PO-07 outright would require, at
minimum, discharging the open-scope items above; the ordering
constraints among those items and the rest of the ledger are taken up
in Chapter~12.
